% discussion.tex -- Section VI (writer W1, 2026-10-30). Binding contract: paper2-tnnls/SPEC.md v2.2, section 3 row VI, sections 7, 8.
% VI-A opens with the A4.2.3 NEG "6.1 opens" sentence; the Tier-3 paragraph and the family H paragraph come after it (A4.1.10 placement
% 2; A5.9 placement 2: own paragraph headed with the prefix, no sentence about P). VI-B is Study 2 alone; the alternatives are the red
% team's items 1-5 (REDTEAM-SPEC "Cheapest alternative explanations"), worded as open questions; item 6 is withdrawn (V2-1) and not used.
% VI-C follows SPEC section 7 item by item (no sentence about the P discordance, V2-1); VI-D follows SPEC section 8.
% No sentence links the two studies; no comparison word; no theory word for a result; every number has a run comment.

\section{Discussion}\label{sec:discussion}

\subsection{Study 1}\label{sec:disc-study1}
% run: runs/P2-GATE-M4/analysis.json:outcome.label (= negative); primary.per_corruption.*.ci90
The preregistered test excludes an improvement or a loss of more than five percentage points in image-direction verification accuracy
from gating under JPEG quality~10, Gaussian noise with $\sigma=0.1$ and a 128-pixel splice, for this model and these severities; it does
not show that gating has no effect.
It concerns this gate at its preregistered level $\alpha=0.05$, one severity per corruption, one model and one data set, and says
nothing about other levels, stronger corruptions, models or tasks, none of which was tested.

% run: runs/P2-GATE-M4/tier3_fusion.json:cells.*.contrasts.* (exploratory; Section IV-C)
The exploratory Tier-3 comparison of the hard gate with five combiners of the same evidence is reported in \cref{sec:study1-results} and
supplement S-II; it describes the combiners and is not an account of the primary result.

% run: runs/P2-GATE-H3/family.json:{n_headroom_items (= 50), n_rule.h_rule (= H-NOGO), n_rule.G_max_headroom (= null), go (= false)}
\emph{In a separately labelled family added after the primary result was known (family H: SugarCrepe hard
negatives~\cite{hsieh2023sugarcrepe} on 50 COCO images not used elsewhere; Holm over the same three corruptions),} the pre-written pilot
rule found no corruption with decision headroom on 50 pilot images, so the family's test was not run, and nothing is stated about the
gate on hard negatives (\cref{sec:study1-H}; supplement S-III).

\subsection{Study 2}\label{sec:disc-study2}
% run: runs/P2-LRF-L1/a6_analysis.json:family_L.L1.label (= confirmed); branch.branch (= A)
L1 shows that one frozen weighting, hash-pinned before any test image was chosen, raised image-direction verification accuracy under
Gaussian noise with $\sigma=0.1$ on 450 COCO images that played no part in its design or fit, passing an exact paired test at the
preregistered level and holding with length-matched negative captions.
It does not show that the weighting works on a corruption it was not trained on, that its weights follow the reliability of a token,
or that the effect carries to another model, data set or severity.
% run: runs/P2-LRF-LG/a6_analysis.json:family_LG.cells.impulse_noise3.ci95 (= [0.0, 0.1]); runs/P2-LRF-L1/a6_analysis.json:family_L.L4.p_holm (= 0.0486)
Five inexpensive alternative readings remain open, each with its evidence in \cref{sec:limitations}: (i) a reweighting that helps
whatever the condition; (ii) a supervised fit to the training corruption; (iii) a single noise-sensitive feature, such as the
codebook-neighbour distance or the one-pass observed-position score, which \cref{fig:lrf_weights} cannot separate and the exploratory
ablation of \cref{sec:study2-ablation} bears on only as an exploratory result; (iv) a limited reach to other noise types, since the
impulse-noise interval starts at 0.0; and (v) an unstable splice decision, whose Holm-adjusted $p$ of 0.049 lies 0.001 below the level.
A second backbone, a corruption held out of training and a preregistered comparison with simpler weighting rules on fresh images are the
tests that would settle these questions; none was run.

\subsection{Limitations}\label{sec:limitations}
\emph{Study 1.} One model and one severity per corruption were tested.

% run: runs/P2-LRF-L1/a6_analysis.json:family_L.L4.p_holm (= 0.0486); runs/P2-LRF-LG/a6_analysis.json:family_LG.cells.impulse_noise3.p_two (= 0.058)
% run: runs/P2-LRF-L1.try1/selfcheck.json:image_tokens_frac_equal (= 0.98828125); runs/P2-LRF-L1.try2/selfcheck.json:image_tokens_frac_equal (= 0.9892578125)
% run: runs/P2-LRF-BASELINES-PARITY-LG/baselines.json:pixel_parity.min (= 0.986328125)
\emph{Study 2.} One model, one data set and one severity per corruption were tested.
LRF was trained on the tested corruption (Gaussian noise with $\sigma=0.1$ was a training condition), so L1 tests new images, not a new
corruption.
No preregistered test separates corruption-specific weighting from a generic gain, since the clean-image and splice estimates were also
positive; the title's ``under Gaussian noise'' names the tested condition, not a specificity.
The Holm-adjusted $p$ of L4 is 0.049, and the two-sided $p$ of the impulse-noise cell is 0.058.
The clean-image margin of L-NI was fixed after an exploratory clean interval had been seen.
Whether COCO images were in the pretraining data of UniDisc, its tokeniser, CLIP or SigLIP is unknown.
Each image has one corruption realisation and one mask seed, so the intervals cover item sampling only.
A fresh tokenisation of the same pixels in another process need not reproduce the stored codes: it agreed on 98.8\,\% and 98.9\,\% of
image tokens in the two stopped self-checks of card row A-8, and on at least 98.6\,\% in the pixel-parity runs.
The reference encoders almost solve this matching task (\cref{sec:study2-reference}); this paper makes no claim of a verifier that ranks
with them.
No ablation was preregistered; \cref{sec:study2-ablation} is exploratory.

\emph{Both studies.} All results are for random-caption negatives. The preregistration is self-hosted in the sense of
\cref{sec:setting-tests}, and the audits were agent-run (see the Declaration).

\subsection{Deviations From the Plans}\label{sec:deviations}
{\small
\noindent\emph{Study 1:}
(1) The Tier-3 comparators, parameters and fallback seed were fixed in an addendum committed before the main run's headroom decision;
their intervals carry no $p$-value and no multiplicity adjustment.
(2) Outcome-label wording was fixed before the primary result, with a dated record of the null wording and of the reading adopted where
card and plan differed on a rejection against the gate.
(3) The time-input statement rests on a code reading at a pinned commit and a key check; its table moved to the supplement.
(4) Family~B ran before the re-infilling family under the plan's fixed time inequality.
(5) Family~H has its own disclosure and amendment items; one decision row was committed after its pilot driver started and before any
pilot row existed.
(6) The effective-$n$ record of the re-infilling figure carries a reason text written at analysis time.
(7) The data steward did not collect the image ids of negatives and mismatches separately (no analysis uses them), and its check left
some run files outside the excluded classes unopened.
(8) The exact Clopper--Pearson bounds of \cref{tab:familyP} were added at writing (\nolinkurl{runs/P2-PAPER-BOUNDS}, descriptive).
\emph{Study 2:}
(1) The L1 stage stopped twice at its tokenisation self-check before any item; card row A-8 moved the code pin, and row A-9 changed one
test only.
(2) The scorer's pin file was written while L1 ran, with every value copied from records committed before L1.
(3) The environment lock files of the design, re-score, L1 and other-noise stages are empty; the environment was recorded after the runs,
and the corruption package's commit was verified from it.
(4) Item arrays are untracked with committed SHA-256 manifests; hashes appear in CRLF and LF forms of the same content.
(5) The pixel-parity files were moved unchanged from a wrong output directory with no prior hash (audit-time hashes are recorded); they
supersede the field ``pixel\_parity: not run'' of the main reference-encoder file.
(6) The reference-encoder download was larger than the size quoted to the first author.
(7) What had been seen before the re-freeze is stated in \cref{sec:study2-provenance}.
(8) ``No output was read before its run was complete'' rests on the operations log.
(9) The exploratory ablation was added at writing.
Supplement S-I gives each item in full with its record.\par}
